% Generated draft: requires mathematical and source review.
\documentclass[11pt]{article}
\usepackage[margin=25mm]{geometry}
\usepackage{fontspec}
\setmainfont{Latin Modern Roman}
\usepackage{amsmath,amssymb,mathtools,unicode-math}
\setmathfont{Latin Modern Math}
\usepackage{enumitem,needspace,xurl,hyperref}
\setlength{\parindent}{0pt}
\setlength{\parskip}{5pt}
\setlength{\emergencystretch}{3em}
\newcommand{\N}{\mathbb{N}}
\newcommand{\Z}{\mathbb{Z}}
\newcommand{\Q}{\mathbb{Q}}
\newcommand{\R}{\mathbb{R}}
\newcommand{\C}{\mathbb{C}}
\newcommand{\F}{\mathbb{F}}
\newcommand{\A}{\mathbb{A}}
\newcommand{\PP}{\mathbb{P}}
\newcommand{\E}{\mathbb{E}}
\newcommand{\eps}{\varepsilon}
\newcommand{\dd}{\,\mathrm{d}}
\begin{document}
\section*{UnsolvedMath ID: 30003244}
\textbf{Source problem code:} catalogue entry
\subsection{Definitions and mathematical statement}
The supplied question is whether the ring \(\mathcal O_{\mathbb Q_{\mathrm{tr}}^{(d)}}\) has Property \((N)\). The record identifies \(\mathbb Q_{\mathrm{tr}}^{(d)}\) as the maximal totally real subfield of \(\mathbb Q^{(d)}\), but does not define \(\mathbb Q^{(d)}\), specify the range or quantification of \(d\), or define Property \((N)\). Consequently, the following explicit formulation is conditional on conventions that require source verification.

Fix an algebraic closure \(\overline{\mathbb Q}\subset\mathbb C\). An algebraic number is totally real if every one of its conjugates over \(\mathbb Q\) lies in \(\mathbb R\). Write \(\mathbb Q_{\mathrm{tr}}\) for the field of all totally real algebraic numbers. Under the conventional interpretation that, for a fixed positive integer \(d\), \(\mathbb Q^{(d)}\) is the compositum in \(\overline{\mathbb Q}\) of all number fields of degree at most \(d\), set
\[
K_d=\mathbb Q^{(d)}\cap\mathbb Q_{\mathrm{tr}},\qquad
\mathcal O_{K_d}=\{\alpha\in K_d:\alpha\text{ satisfies a monic polynomial in }\mathbb Z[x]\}.
\]
Thus \(K_d\) is the maximal totally real subfield described in the record.

For an algebraic number \(\alpha\) with primitive irreducible polynomial \(a\prod_{i=1}^{n}(x-\alpha_i)\in\mathbb Z[x]\), where \(a>0\), define its absolute multiplicative Weil height by
\[
H(\alpha)=\left(a\prod_{i=1}^{n}\max\{1,|\alpha_i|\}\right)^{1/n}.
\]
If Property \((N)\) means the Northcott property for this height, the question becomes: does
\[
\forall B\in\mathbb R_{>0},\qquad
\#\{\alpha\in\mathcal O_{K_d}:H(\alpha)\le B\}<\infty
\]
hold for the intended values of \(d\)? No bound on the degrees of the elements \(\alpha\) is imposed in this formulation.

A different possible interpretation uses the house \(\max_{1\le i\le n}|\alpha_i|\). Under that interpretation the question instead asks whether
\[
\forall B\in\mathbb R_{>0},\qquad
\#\left\{\alpha\in\mathcal O_{K_d}:\max_{1\le i\le n}|\alpha_i|\le B\right\}<\infty.
\]
The supplied record does not specify which boundedness condition defines Property \((N)\); these formulations are not identified with one another here. The intended definition and quantification of \(d\) remain unresolved.
\subsection{Short English statement}
The source asks whether the algebraic integers in the maximal totally real subfield of \(\mathbb Q^{(d)}\) have Property \((N)\). A precise interpretation requires checking the definition of \(\mathbb Q^{(d)}\), the intended values of \(d\), and whether Property \((N)\) concerns bounded Weil height or bounded house. The supplied open status is dated curation, not an independently verified conclusion.
\subsection{Sources}
\begin{itemize}
\item Supplied local registry record: permanent UnsolvedMath ID \(30003244\), title ``Northcott Property for Totally Real Integer Rings,'' source problem code \texttt{catalogue entry}, in the Oberwolfach Reports Open Problems collection. The record gives proposed year 2016 and names no individual proposer.
\item Supplied report citation: \emph{Definability and Decidability Problems in Number Theory} (2017), Oberwolfach Reports, pp.\ 2793--2866, DOI \texttt{10.4171/owr/2016/49}. \url{https://doi.org/10.4171/owr/2016/49}. The background describes these reports as documenting mathematical workshops and their open-problem discussions; no precise location within the report is supplied for this question.
\item Supplied literature source, labelled ``Oberwolfach Reports 2017'': \url{https://ems.press/content/serial-article-files/46652}.
\item Supplied literature source, labelled ``Daans et al., Universal quadratic forms'': \url{https://doi.org/10.1112/jlms.70022}. No further bibliographic details or supporting passage are supplied.
\item Literature assessment dated exactly \texttt{2026-08-21}: the supplied curation labels the question open and states that the totally real integer ring is still treated as an open Northcott-property question. This is supplied dated curation derived from the OWR record, not independently verified literature assessment; the record itself requests independent verification of time-sensitive status and cited claims.
\item The registry labels the statement \texttt{corrected\_verified} and reports a prior manual semantic, typographic, and LaTeX review claiming that the clean statement is self-contained and preserves its target. These are supplied registry assertions, not mathematical verification of this draft or resolution of its missing definitions.
\end{itemize}
\end{document}
